\section{Limitations}
\label{sec:limitations}

We collect the limitations of the design under five headings, each
of which suggests a corresponding direction for follow-up work.

\subsection{Single news feed and proxy timestamps}

The analysis uses one news feed (the FinancialJuice Discord
newsfeed) and the Discord publication timestamp as the event time.
The publication timestamp is downstream of the underlying primary
source---typically a wire service, an agency feed, or an official
release calendar---and the variable latency between the primary
source and the Discord publication is, in our reading, the most
plausible explanation for the pre-event drift reported in
Section~\ref{ssec:results-timing}. We cannot decompose this drift
into ``feed latency'' and ``true pre-disclosure activity'' with the
data at hand. Future work should either obtain primary-source
timestamps directly, or compare a sample of events against a
higher-tier feed (e.g., Bloomberg, Reuters Newscope) to estimate
the latency distribution.

A consequence of using a single feed is that any feed-specific
bias in the curation of ``gold'' events propagates into the sample.
The gold filter is deterministic and reproducible, but it relies
on the publisher's own red-dot and \texttt{\$MACRO} tagging
conventions. We do not have a counterfactual sample drawn from a
different feed.

\subsection{Two-asset universe and CFD proxy for NDX}

We test two assets, EUR/USD spot and the Nasdaq-100 CFD. The
choice reflects the most liquid headline-FX pair and a tradable
index proxy. The findings should not be extrapolated to
less-liquid currency pairs, to individual equities (where
firm-specific news effects dominate), or to fixed-income markets
without explicit testing. The Dukascopy \texttt{E\_NQ-100} series
is a CFD price track for the Nasdaq-100, not the authoritative
consolidated tape; we use it because it offers one-minute
resolution with no gaps over our period, but readers should treat
the magnitude estimates on NDX as estimates on the CFD price
rather than on the cash index.

\subsection{Volume series is a proxy}

The volume series accompanying the Dukascopy NDX CFD is best
understood as tick volume (the count of price updates) rather than
consolidated exchange volume. The H10 test of event-window volume
elevation is reported in the Online Appendix with this caveat
explicit; readers should not interpret it as evidence on changes
in consolidated traded volume around news events. A defensible
test of the volume question requires consolidated tape access for
the index constituents.

\subsection{Sentiment classifier not yet fully manually validated}

We benchmark the LLM-derived sentiment against two reference
classifiers (Loughran-McDonald dictionary, FinBERT-tone). We have
prepared a 200-event subsample for double-coded human annotation
with the infrastructure to compute Cohen's $\kappa$ between
annotators, $F_1$ for the LLM against the consensus label, and a
split-by-knowledge-cutoff comparison to address memorisation
concerns. At the time of writing the annotation is in progress;
we update the paper with the results when complete. The
pre-cutoff/post-cutoff stability of the central magnitude findings
(Section~\ref{ssec:results-robustness}, C5) provides indirect
evidence against the memorisation hypothesis, but is not a
substitute for direct human validation of the labels.

Additionally, the LLM-reported \texttt{confidence} field is
uncalibrated (H6, reported here as a null finding): the model's
stated probability that its direction is correct does not
correspond to the empirical hit rate within confidence buckets. We
do not use the confidence as a probability anywhere in the
analysis, but its inclusion in the LLM output schema invites
misuse and should be removed or re-calibrated in future versions
of the pipeline.

Two further null findings should be reported explicitly. H12 (a
test of bear-versus-bull asymmetry in the price response) does not
detect a robust asymmetry in our sample; this is a limitation of
either the sample, the test, or the underlying phenomenon, and we
cannot distinguish among these. C8 (a cross-model consistency
check on $n=200$ events comparing \texttt{deepseek-v4-flash} and
\texttt{deepseek-v4-pro}) is exploratory: the subsample is
underpowered for the inference suggested by the headline numbers,
and we report it in the Online Appendix rather than as a main
result.

\subsection{Multiple testing, dependence, and pre-specification}

The pipeline runs more than a hundred individual $p$-values across
the main tests, the auxiliary tests, and the cells of stratified
analyses. We address this with uniform Benjamini-Hochberg false
discovery rate adjustment \citep{benjamini1995controlling}, and we
draw the primary conclusions exclusively from $q$-values. Two
residual concerns remain.

First, the cluster-dependence structure means that adjacent events
in a news burst are not independent. We address this with the
cluster-dedupe convention in non-regression tests and with
cluster-robust standard errors in regression tests; both devices
may under-correct in long news bursts (cluster sizes occasionally
exceed twenty events).

Second, the C1--C7 extensions were added during methodological
revision after the H1--H14 pre-specification. The C extensions are
clearly demarcated in the code and in \texttt{STATUS.md}. Readers
who wish to apply a stricter pre-registration standard can re-read
Section~\ref{sec:results} attending only to the H1--H14 results;
the central conclusions of
Section~\ref{ssec:results-volatility}~(H1) and
Section~\ref{ssec:results-timing}~(H8, H9) survive this
restriction.

\subsection{What would strengthen a follow-up}

In addition to the manual validation and the primary-source
timestamp work mentioned above, the most promising extensions are:
(a)~cross-feed replication on a second public news service to
quantify feed-specific bias; (b)~cross-asset extension to a small
basket of individual equities for firm-specific news; (c)~prompt
engineering and LLM ensembling on the directional task to test
whether the modest direction edge can be improved beyond the
present $2$--$6$ pp margin over baselines; and (d)~integration of
the methodology with an order-flow or quote-imbalance dataset to
distinguish the magnitude response into a re-pricing component and
a liquidity-withdrawal component.
